% F-19 prototype: single shared preamble for all chapter bodies.
% Union of packages/macros/theorem environments found across chap01-13
% (see INVENTARIO.md). Macros use \providecommand so the FIRST definition
% wins; the handful of cosmetic variants found by inventory.py
% (\mathrm{BV} vs \operatorname{BV}, etc.) render identically.

% morewrites MUST load before hyperref/other packages claim \write streams:
% 13 chapters x (bibliography + hyperref + float aux) exhaust the classic
% 16-stream limit; morewrites raises the ceiling to 9000+ via etex registers.
\usepackage{morewrites}
\usepackage[utf8]{inputenc}
\usepackage[T1]{fontenc}
\usepackage{amsmath,amssymb,amsfonts,amsthm,mathtools}
\usepackage{mathrsfs}
\usepackage{microtype}
\usepackage{booktabs}
\usepackage{tabularx}
\usepackage{array}
\usepackage[dvipsnames]{xcolor}
\usepackage{hyperref}
\usepackage{geometry}
\usepackage{enumitem}
\usepackage{graphicx}
% chap13's \includegraphics uses a bare filename (figure lives next to the
% chapter source, one level up from this prototype folder); point graphicx
% there instead of editing the chapter body.
\graphicspath{{../}}
\usepackage{cite}
\usepackage{silence}
\WarningFilter{nameref}{The definition of \label has changed!}
\usepackage{tikz}
% chap13 is revtex4-2; extract.py rewrites its \ruledtabular table into a
% plain booktabs \toprule/\bottomrule + center (see extract.py), so no
% \ruledtabular shim is needed here.
% revtex4-2's \acknowledgments environment (chap13) is just a run-in
% "Acknowledgments." heading; book has no built-in equivalent.
\newenvironment{acknowledgments}
  {\par\medskip\noindent\textbf{Acknowledgments.}\ }
  {\par}

\numberwithin{equation}{chapter}

% Paragraph reflow: each chapter was authored against its OWN margin
% (0.9in for ch. 3-6, 1in for ch. 7-9, a custom \geometry for ch. 1-2,
% class defaults for ch. 10-13) before being unified here to 3cm/2.8cm
% (the real master's own geometry). Raising tolerance/emergency stretch
% is the standard, honest way to let TeX re-justify those paragraphs at
% the new width instead of flagging every minor over/underfull line;
% it changes glue stretch, not the visible word choice.
\tolerance=1000
\emergencystretch=2em

% F-19 adoption: attempt to clear the remaining overfull \hbox WITHOUT
% editing any chapter. \allowdisplaybreaks lets long multi-line align/
% gather blocks break across pages instead of being forced into a cramped
% single page. A first attempt also tried \everydisplay{\small} (a
% smaller display-math font); REVERTED -- \small is a text-mode size
% switch and \everydisplay inserts it already inside math mode, which
% raised a "Command \small invalid in math mode" LaTeX Font Warning on
% every single display equation in the book (dozens), failing the
% 0-warnings gate outright, and even made the worst remaining overfull
% box bigger (64pt vs. 26pt) through inconsistent math sizing. Any
% working display-font-shrink would need a different hook (e.g.
% \AtBeginEnvironment{equation}{\small} from etoolbox, repeated for every
% display environment the chapters use) and is left for a future pass;
% see ADOCAO_F19.md for the resulting overfull count with just
% \allowdisplaybreaks plus the existing \tolerance/\emergencystretch.
\allowdisplaybreaks[4]

% --- Theorem environments (union across chapters; numbered per chapter) ---
\theoremstyle{plain}
\newtheorem{theorem}{Theorem}[chapter]
\newtheorem{lemma}[theorem]{Lemma}
\newtheorem{proposition}[theorem]{Proposition}
\newtheorem{corollary}[theorem]{Corollary}
\newtheorem{conjecture}[theorem]{Conjecture}
\newtheorem{problem}[theorem]{Problem}
\newtheorem{algorithm}[theorem]{Algorithm}

\theoremstyle{definition}
\newtheorem{definition}[theorem]{Definition}
\newtheorem{example}[theorem]{Example}
\newtheorem{construction}[theorem]{Construction}
\newtheorem{assumption}[theorem]{Assumption}

\theoremstyle{remark}
\newtheorem{remark}[theorem]{Remark}

% --- Custom notation: AUTO-HARVESTED from every chapter's preamble by
% extract.py (union, first definition wins, each as \providecommand).
% See gen/common_macros.tex. Do not hand-list macros here: re-run
% extract.py after touching a chapter's \newcommand/\providecommand block.
\input{gen/common_macros}

% 2026-10-11: \xref{label}{fallback} is harvested above as a standalone
% fallback (\providecommand{\xref}[2]{#2}, defined in whichever chapter's
% preamble extract.py scans first -- see chap05/09/11/12). In the master,
% EVERY chapter shares one numbering, so override it here to resolve the
% real cross-chapter number via \ref; #1 is already the target label with
% its chapNN: prefix (e.g. chap07:thm:lower_bounds), untouched by
% extract.py's own \ref-prefixing because \xref is not in its REF_CMDS list.
\renewcommand{\xref}[2]{\ref{#1}}

% book's default "headings" pagestyle puts the CURRENT SECTION title
% (via \sectionmark/\rightmark) in the running head on right-hand pages.
% Several chapters have section titles far longer than the text width
% (e.g. "2.6. TENSOR NETWORKS: GAUGE THEORY, PATH INTEGRALS, AND CONTINUOUS
% LIMITS"), which overflows the head rule regardless of \chapter[short]{}.
% Drop section titles from the running head; keep the chapter's short
% title (set via \chaptermark in the generated master) on both sides.
\renewcommand{\sectionmark}[1]{}

% --- Manual "Abstract" block (book class has no native abstract env;
% \quotation/\endquotation is a standard LaTeX kernel environment) ---
\newenvironment{chapterabstract}
  {\begin{quotation}\noindent\textbf{Abstract.}\ }
  {\end{quotation}}
